\begin{abstract}
Bank support desks must choose which ticket to handle next, as urgent cases such as fraud have a limited response window. \af{Too many things cluttered in the opening sentence. Rephrase to clarify the problem}Treating ticket prioritization as a classification task, we train multilingual models on an offline labelled dataset to predict each ticket's intent, sentiment, and priority. Fine-tuned LaBSE achieves the best pooled intent (0.8835 macro-F1); sentiment (0.7138 Negative-F1) and priority (0.8901 macro-F1) it ties with other encoders, small decoders, and the SVM within their confidence intervals. \af{this as a classification problem which is trained with an offline dataset....and add how the rest is handled interms of}\af{next sentence not clear...rephrase. Sentences do not start with because}Intent accounts for 76.9\% of the uncertainty in priority, treating classification heads as independent evidence counts shared signal twice; an equal-weight log pool merges direct priority predictions with an intent-conditioned chain, reducing log loss by 27.4\% and resolving collinearity. We then specify a dynamic urgency scoring function that couples this combined posterior with elapsed waiting time, so that queue position reflects both predicted severity and accrued delay and no ticket can be starved indefinitely by construction. \af{dont add future work in abstract}To train and evaluate the system, we expanded BANKING77 into a corpus of 65,385 tickets across five language variants, pairing LLM-generated priority and sentiment labels, with sentiment-label quality assessed against a manually annotated gold subset.
\end{abstract}